\par\noindent\begin{minipage}{\linewidth}
% !TeX root = ../教材题库_全部条目_转写工作稿.tex
% Source: XJ-B1-C02-EX-P057-Y0534
% 第2章 / 2.3 一元二次不等式 / 教材 p.57 / PDF 第64页
% 已对照原页转写题干及公式；待教师复核；未制作答案
% 原题及相邻原解截图（核对用）：../题目原图/XJ-B1-C02-EX-P057-Y0534-01.webp
% 本文件只在转写完成后作为题干使用；OCR 条目尚不是可靠题干。
\begin{bankitem}{2.3 例题 例题9}{第2章 / 2.3 一元二次不等式 / 教材 p.57 / PDF 第64页}{XJ-B1-C02-EX-P057-Y0534}{已对照原页转写题干及公式；待教师复核；未制作答案}
某化学试剂生产厂以 $x\,\mathrm{kg/h}$ 的速度运输生产某种产品（生产条件要求边生产边运输，且 $1\le x\le10$），每小时可获得利润 $100\left(5x+1-\dfrac3x\right)$ 元。
\subq{1}要使运输生产该产品 $2\,\mathrm h$ 获得的利润不低于 3000 元，求 $x$ 的取值范围。
\subq{2}要使运输生产 $900\,\mathrm{kg}$ 该产品获得的利润最大，该工厂应该选取何种运输生产速度？并求最大利润。
\end{bankitem}
\end{minipage}\par\medskip
